\documentclass[journal, a4paper]{IEEEtran}
\usepackage{booktabs}
\usepackage{graphicx}
\usepackage{stfloats}
\usepackage{placeins}
\usepackage{url}
\usepackage{cite}
\usepackage{cuted}
\usepackage{amssymb}
\usepackage{lipsum}
\usepackage{amsmath}
\usepackage{siunitx}
\sisetup{output-exponent-marker=\ensuremath{\mathrm{e}}}

\begin{document}

\title{Allocating Model-Based Contribution to Market Externalities: Coalition Shapley Values, Local Participation Sensitivities, and a Pigouvian Benchmark}

\author{Jan~Ol\v{s}ina\\[1ex] Independent Researcher, Prague, Czech Republic\\
email: \texttt{jan.olsina@gmail.com},  ORCID: \texttt{0009-0003-9816-3958}.}


\maketitle

\begin{abstract}
When a competitive market outcome generates a harmful externality (e.g., CO$_2$ emissions), an analyst may wish to allocate the modeled harm among heterogeneous producers and consumers. This paper compares three curve-based allocation rules. First, a cooperative-game rule applies the Shapley value to a cost game whose characteristic function is the harm generated by coalition-specific market clearing. Second, a local participation-sensitivity rule uses comparative statics: each agent's weight is the sensitivity of equilibrium quantity to a small proportional scaling of that agent's supply or demand curve. The resulting producer--consumer split can be written in terms of local slopes or elasticities. For nonlinear harm, converting these local quantity sensitivities into a budget-balanced allocation of the harm level is an average-damage normalization, not an exact differential decomposition. Third, a Pigouvian benchmark separates harm at the welfare-optimal quantity from excess harm relative to that benchmark; a further beneficiary-pays rule allocates the excess in proportion to private surplus gains from the absence of the corrective tax. The Shapley formula, comparative statics, and Pigouvian benchmark are established results. The paper's proposed contribution is their joint application to one producer--consumer attribution problem, the particular coalition market-clearing cost game, the closed-form symmetric linear example and limits, and the explicit local elasticity formula. These are model-based attribution rules, not by themselves measures of actual causation, legal liability, or moral blame.
\end{abstract}

\section{Introduction}
Markets coordinate production and consumption through prices, but when traded goods generate harmful externalities, the equilibrium quantity can exceed the socially desirable level. Beyond standard policy prescriptions (taxes, caps, regulation), there is a normative question: \emph{how should a modeled harm total be attributed between producers and consumers}, especially when each agent has a different supply or demand curve?

This paper formalizes three complementary allocation rules based on primitives familiar from microeconomic theory. The first is a coalition attribution: the Shapley value is applied to a cost game generated by coalition-specific competitive equilibria. The second is a local participation-sensitivity rule: weights are based on each agent's marginal influence on equilibrium quantity under a small proportional perturbation of that agent's curve. The third is a counterfactual rule: the unregulated equilibrium is compared with a Pigouvian benchmark, and the resulting excess-harm gap is allocated according to private benefits from the absence of corrective pricing.

The underlying ingredients are mostly standard. The Shapley value and its random-order marginal-contribution interpretation originate with Shapley~\cite{Shapley1953}; continuous or infinitesimal allocation by path-integrated marginal contributions is associated with Aumann and Shapley~\cite{AumannShapley1974}. The distinction between private and social net product and the corrective-tax idea go back to Pigou~\cite{Pigou1932}, with later formal treatments and qualifications in, among others, Baumol~\cite{Baumol1972} and Coase~\cite{Coase1960}. Local comparative statics and tax-incidence reasoning are also standard microeconomics. The proposed contribution here is narrower: a particular coalition market-clearing harm game, a closed-form symmetric linear specialization with limiting producer shares, an explicit proportional-participation elasticity allocation, and a side-by-side comparison with a Pigouvian beneficiary-pays benchmark.

A key theme is that any such allocation depends on (i) a harm model, (ii) an equilibrium or market-clearing institution, (iii) a perturbation or coalition convention, and (iv) where relevant, a counterfactual baseline. We keep the environment simple (single good, competitive behavior, quantity-only harm) to distinguish algebraic results from normative choices.


\subsection{Relation to prior work and scope of the novelty claim}
The Shapley formula, its efficiency, symmetry, null-player, and additivity properties, and its interpretation as an average marginal contribution over random orderings are established results~\cite{Shapley1953}. Shapley-based allocation has long been used for joint costs and other attribution problems, while Aumann--Shapley prices provide a continuous path-based analogue~\cite{AumannShapley1974,BilleraHeath1982}. Cooperative market games are also classical~\cite{ShapleyShubik1969}, although their coalition worth is ordinarily defined from attainable utility rather than from external harm at an internally cleared coalition market. Shapley methods are well established in sensitivity analysis~\cite{Owen2014,SongNelsonStaum2016} and have been used for carbon or energy-system allocation problems~\cite{XieChen2023,FengEtAl2025}. Therefore, neither ``using Shapley values for harm'' nor ``using marginal sensitivities for attribution'' is, by itself, novel.

The specific claims of contribution made in this paper are instead:
\begin{enumerate}
\item the characteristic function $v(T)=H(Q(T))$ obtained by clearing a competitive market separately for every producer--consumer coalition under the stated no-trade convention;
\item the finite-sum expression and limiting producer-side Shapley shares for identical truncated-linear agents in Section~\ref{sec:shapley_linear_identical};
\item the local proportional-participation sensitivity formulas in Section~\ref{sec:elasticity}, including the cross-side elasticity split and the explicit warning that the budget-balanced harm allocation is a normalization when $H$ is nonlinear; and
\item the comparison of those two realized-outcome allocations with a Pigouvian benchmark and a separately identified beneficiary-pays rule.
\end{enumerate}
These claims should be read as claims of formulation and synthesis, not as a claim that no mathematically equivalent construction exists elsewhere. Establishing priority would require a more exhaustive literature review than is attempted here.

\subsection{Attribution versus Causation and Moral Blame}
\label{sec:contribution_vs_blame}
Before proceeding, it is worth being precise about what the quantities defined in this paper do and do not represent. The first two frameworks decompose a modeled harm outcome into agent-level shares, while the third combines a welfare counterfactual with an explicitly normative beneficiary-pays rule. The Shapley construction is an allocation in a specified cooperative cost game; the local construction is an intervention-like sensitivity to a scaling parameter; and the Pigouvian construction is a benchmark comparison. None is automatically a measure of actual causation in the structural-causal sense. In particular, a Shapley value becomes ``causal'' only if its coalition operation is defended as an appropriate intervention, a distinction emphasized in the causal-attribution literature~\cite{HalpernPearl2005,HeskesEtAl2020}.

\emph{Moral blame} (or culpability, fault, responsibility in the everyday normative sense) is a different and stricter notion. Two agents with the same modeled physical contribution to harm are not automatically equally blameworthy: blame typically also depends on whether the agent could foreseeably have acted otherwise, whether a less harmful alternative was reasonably available, whether the agent possessed relevant information, and whether the conduct fell short of a standard of care. A producer using the best available technology and one using needlessly dirty technology may generate the same emissions, yet few moral frameworks would treat them as equally at fault. Conversely, an agent assigned a large model-based share may attract little blame if realistic alternatives were unavailable.

Because the constructions in this paper are model-dependent allocations, we generally use \emph{attribution}, \emph{contribution weight}, or \emph{sensitivity} rather than \emph{blame} or \emph{culpability}. Formal accounts of responsibility and blame require additional causal and epistemic structure~\cite{ChocklerHalpern2004}. Section~\ref{sec:exante} therefore treats conditioning on information only as an expected-attribution exercise, not as a complete theory of culpability.

\section{Setup}
\subsection{Agents and curves}
There are $N$ producers indexed by $i\in \mathcal{P}=\{1,\dots,N\}$ and $M$ consumers indexed by $j\in \mathcal{C}=\{1,\dots,M\}$. Let the set of all agents be $\mathcal{N}=\mathcal{P}\cup \mathcal{C}$.

\paragraph{Supply.}
Producer $i$ supplies
\begin{equation}
s_i(p)\ge 0,\qquad s_i'(p)\ge 0,
\end{equation}
for price $p\ge 0$.

\paragraph{Demand.}
Consumer $j$ demands
\begin{equation}
d_j(p)\ge 0,\qquad d_j'(p)\le 0.
\end{equation}

\paragraph{Aggregation.}
Aggregate supply and demand are
\begin{equation}
S(p)=\sum_{i\in \mathcal{P}} s_i(p),\qquad D(p)=\sum_{j\in \mathcal{C}} d_j(p).
\end{equation}

\subsection{External harm}
We assume a \emph{quantity-only} externality: total external harm depends on total traded quantity $Q$ via
\begin{equation}
H(Q),\qquad H(0)=0,\qquad H'(Q)\ge 0.
\end{equation}
We interpret $H(Q)$ as monetized damages, or any agreed scalar measure of harm. The marginal harm is
\begin{equation}
MH(Q) \triangleq H'(Q).
\end{equation}

\subsection{Competitive equilibrium}
A (Walrasian) competitive equilibrium price $p^\star$ solves
\begin{equation}
S(p^\star)=D(p^\star),
\end{equation}
with equilibrium quantity
\begin{equation}
Q^\star = S(p^\star)=D(p^\star),
\end{equation}
and realized harm
\begin{equation}
H^\star \triangleq H(Q^\star).
\end{equation}

\section{Shapley Attribution for Realized Harm}
\subsection{Characteristic function via coalition markets}
To allocate the modeled harm, we define a cooperative cost game in which the worth of a coalition is the harm generated when only members of that coalition participate in the market.

Let $T\subseteq \mathcal{N}$ be a coalition. Define coalition supply and demand:
\begin{equation}
S_T(p)=\sum_{i\in T\cap \mathcal{P}} s_i(p),\qquad
D_T(p)=\sum_{j\in T\cap \mathcal{C}} d_j(p).
\end{equation}
A coalition market-clearing price $p(T)$ solves
\begin{equation}
S_T(p(T))=D_T(p(T)),
\end{equation}
whenever $T$ contains at least one producer and one consumer and a clearing price exists. We assume that the resulting traded quantity is unique; alternatively, a fixed tie-breaking rule must be specified. Define
\begin{equation}
Q(T)=S_T(p(T))=D_T(p(T)).
\end{equation}
For coalitions that cannot trade (no producers or no consumers, or no clearing), we set $Q(T)=0$. This no-trade convention is part of the proposed model, not a property of the Shapley value. Different rationing, outside-market, or disposal rules generally produce a different game and a different allocation. Outsiders are removed rather than allowed to form a separate market, so $v$ is a characteristic-function game. If coalition worth depended on how outsiders were partitioned, a partition-function formulation would be required~\cite{ThrallLucas1963}. The economic external harm in $H$ should not be confused with coalition externalities in that game-theoretic sense.

\paragraph{Quantity-only harm game.}
Define the characteristic function
\begin{equation}
v(T)\triangleq H(Q(T)).
\label{eq:charfunc}
\end{equation}
This encodes the idea that harm is jointly produced by the coalition's capacity to supply and willingness to demand.

\subsection{Shapley value}
The Shapley value~\cite{Shapley1953} assigns each agent $a\in \mathcal{N}$ a cost allocation
\begin{equation}
\phi_a(v)\;=\mspace{-13mu}\sum_{T\subseteq \mathcal{N}\setminus\{a\}}
\frac{|T|!\,(|\mathcal{N}|-|T|-1)!}{|\mathcal{N}|!}\Big(v(T\cup\{a\})-v(T)\Big).
\label{eq:shapley}
\end{equation}
Interpretation: $\phi_a(v)$ is the expected marginal change in the coalition cost when agent $a$ is inserted into a uniformly random ordering of all agents. Calling this change ``causal'' requires the additional substantive assumption that coalition entry is a meaningful intervention.

\paragraph{Budget balance.}
Because $v(\varnothing)=0$, the Shapley value satisfies
\begin{equation}
\sum_{a\in \mathcal{N}} \phi_a(v)=v(\mathcal{N})=H(Q^\star)=H^\star,
\end{equation}
so the allocated costs sum to realized harm. Efficiency alone does not establish that the allocation is morally fair or causally correct.

\subsection{Quantity-only specialization}
A commonly used special case sets
\begin{equation}
H(Q)=\kappa Q,\qquad \kappa>0,
\label{eq:linearharm}
\end{equation}
so that $v(T)=\kappa Q(T)$. Then Shapley contribution is linear in coalition quantities:
\begin{equation}
\phi_a(v)=\kappa\cdot \phi_a(Q),
\end{equation}
where $\phi_a(Q)$ is computed from \eqref{eq:shapley} with $v(T)=Q(T)$.


\subsection{Closed-form Shapley attribution for identical linear agents}
\label{sec:shapley_linear_identical}
This subsection provides an explicit Shapley attribution in a tractable special case: $N_P$ identical producers and $N_C$ identical consumers with linear (truncated) curves, and quantity-only harm. The goal is to make the dependence of the producer--consumer split on slope parameters transparent.

\paragraph{Identical truncated-linear curves.}
Let each producer have supply
\begin{equation}
q^s(p)=\max\{0,\ \alpha p-\gamma\},\qquad \alpha>0,
\label{eq:lin_supply}
\end{equation}
and each consumer have demand
\begin{equation}
q^d(p)=\max\{0,\ \delta-\beta p\},\qquad \beta>0.
\label{eq:lin_demand}
\end{equation}
Define the overlap parameter
\begin{equation}
X \triangleq \alpha\delta-\beta\gamma.
\label{eq:X_def}
\end{equation}
If $X\le 0$ (equivalently $\delta/\beta\le \gamma/\alpha$), there is no price at which both sides are simultaneously positive, so every coalition trades $0$ and all Shapley contributions are $0$. Below we assume $X>0$.

\paragraph{Coalition equilibrium quantity.}
For any coalition containing $m\in\{0,1,\dots,N_P\}$ producers and $n\in\{0,1,\dots,N_C\}$ consumers, aggregate (untruncated) supply and demand are
\begin{equation}
Q_m^s(p)=m(\alpha p-\gamma),\qquad Q_n^d(p)=n(\delta-\beta p).
\end{equation}
For $m,n>0$, coalition market clearing yields
\begin{equation}
p(m,n)=\frac{n\delta+m\gamma}{m\alpha+n\beta},
\end{equation}
and the cleared coalition quantity is
\begin{equation}
Q(m,n)
=\frac{mn(\alpha\delta-\beta\gamma)}{m\alpha+n\beta}
=\frac{mnX}{m\alpha+n\beta}.
\label{eq:Qmn}
\end{equation}
When $X>0$, the clearing price $p(m,n)$ lies between $\gamma/\alpha$ and $\delta/\beta$, so truncation does not bind and \eqref{eq:Qmn} is valid. We therefore define the quantity-only characteristic function
\begin{equation}
v(m,n)=
\begin{cases}
0, & m=0\ \text{or}\ n=0,\\[4pt]
\displaystyle H(Q(m,n)), & m,n>0,
\end{cases}
\end{equation}
and in the linear-harm case \eqref{eq:linearharm} this reduces to $v(m,n)=\kappa Q(m,n)$.

\paragraph{Marginal contribution of an additional producer.}
Fix a particular producer. Consider a predecessor set in a Shapley permutation containing
\begin{align*}
i&\in\{0,\dots,N_P-1\}\ \text{other producers},\\
j&\in\{0,\dots,N_C\}\ \text{consumers}.
\end{align*}
For $j=0$, trade is impossible and the marginal contribution is zero. For $j\ge 1$, the marginal increase in coalition quantity from adding one producer is
\begin{align}
\Delta_P(i,j)
&\triangleq Q(i+1,j)-Q(i,j)\nonumber\\
&=\frac{j^2\,\beta\,X}{(i\alpha+j\beta)\big((i+1)\alpha+j\beta\big)}.
\label{eq:DeltaP}
\end{align}

\paragraph{Exact Shapley value for one producer (finite sum).}
Let $n_{\mathrm{tot}}=N_P+N_C$ be the total number of agents. In the Shapley formula \eqref{eq:shapley}, each predecessor set $T$ with $|T|=k$ carries weight
\begin{equation}
\frac{k!\,(n_{\mathrm{tot}}-k-1)!}{n_{\mathrm{tot}}!}
=\frac{1}{n_{\mathrm{tot}}}\cdot\frac{1}{\binom{n_{\mathrm{tot}}-1}{k}}.
\label{eq:shapley_weight_binom}
\end{equation}
The number of predecessor sets with $(i,j)$ composition is $\binom{N_P-1}{i}\binom{N_C}{j}$, and $k=i+j$. Thus, for linear harm $H(Q)=\kappa Q$ the Shapley allocation to one producer is
\begin{equation}
\phi_P
=
\kappa\sum_{i=0}^{N_P-1}\sum_{j=1}^{N_C}
\binom{N_P-1}{i}\binom{N_C}{j}\;
\frac{1}{n_{\mathrm{tot}}}\frac{1}{\binom{n_{\mathrm{tot}}-1}{i+j}}\;
\Delta_P(i,j),
\label{eq:phiP_general}
\end{equation}
with $\Delta_P(i,j)$ given by \eqref{eq:DeltaP}. By symmetry, all producers receive the same $\phi_P$.

\paragraph{Total producer-side Shapley contribution.}
The combined Shapley allocation to all producers is
\begin{equation}
B_P^{\mathrm{Sh}}=N_P\,\phi_P,
\label{eq:BPSh_def}
\end{equation}
and budget balance implies $B_P^{\mathrm{Sh}}+B_C^{\mathrm{Sh}}=\kappa Q(N_P,N_C)$ where
\begin{equation}
Q(N_P,N_C)=\frac{N_PN_CX}{N_P\alpha+N_C\beta}.
\label{eq:Q_grand}
\end{equation}
It is often convenient to consider the producer share of the realized-harm allocation,
\begin{equation}
s_P^{\mathrm{Sh}}
\triangleq
\frac{B_P^{\mathrm{Sh}}}{\kappa Q(N_P,N_C)}.
\label{eq:sP_def}
\end{equation}
Because both numerator and denominator are proportional to $X$, the share $s_P^{\mathrm{Sh}}$ depends on $\gamma,\delta$ only through the feasibility condition $X>0$, and otherwise depends on $(N_P,N_C)$ and the slope ratio
\begin{equation}
r\triangleq \frac{\alpha}{\beta}.
\label{eq:r_def}
\end{equation}

\paragraph{Simple limiting cases.}
The Shapley share \eqref{eq:sP_def} admits transparent limits as $r$ varies.

\begin{itemize}
\item \textbf{Inelastic supply ($r\to 0$).} In this limit, for $j\ge 1$ the coalition quantity becomes essentially independent of $j$:
\begin{equation}
Q(m,n)=\frac{mnX}{m\alpha+n\beta}\xrightarrow[r\to 0]{}\frac{mX}{\beta}\qquad (n>0),
\end{equation}
so each additional producer adds approximately $X/\beta$ once any consumer is present. Under random arrival orderings, the expected fraction of producers arriving after the first consumer is $N_C/(N_C+1)$, yielding
\begin{equation}
s_P^{\mathrm{Sh}} \xrightarrow[r\to 0]{} \frac{N_C}{N_C+1},
\qquad
1-s_P^{\mathrm{Sh}} \xrightarrow[r\to 0]{} \frac{1}{N_C+1}.
\label{eq:limit_r0}
\end{equation}

\item \textbf{Elastic supply ($r\to \infty$).} In this limit, for $m\ge 1$ the coalition quantity becomes essentially independent of $m$:
\begin{equation}
Q(m,n)\xrightarrow[r\to \infty]{}\frac{nX}{\alpha}\qquad (m>0),
\end{equation}
so only the first producer to arrive matters. The producer-side Shapley share tends to
\begin{equation}
s_P^{\mathrm{Sh}} \xrightarrow[r\to \infty]{} \frac{1}{N_P+1},
\qquad
1-s_P^{\mathrm{Sh}} \xrightarrow[r\to \infty]{} \frac{N_P}{N_P+1}.
\label{eq:limit_rinf}
\end{equation}

\item \textbf{Symmetry point ($N_P=N_C$ and $r=1$).} When the two groups have equal counts and equal slopes ($\alpha=\beta$), the value function $Q(m,n)$ is invariant under exchanging the labels ``producer'' and ``consumer''. By Shapley symmetry, the total allocation must split evenly:
\begin{equation}
N_P=N_C,\ r=1
\quad\Rightarrow\quad
s_P^{\mathrm{Sh}}=\frac{1}{2}.
\label{eq:symmetry_half}
\end{equation}
\end{itemize}

\paragraph{Remark (comparison to local sensitivity).}
The limiting formulas \eqref{eq:limit_r0}--\eqref{eq:limit_rinf} provide a global, order-averaged notion of contribution. By contrast, Section~\ref{sec:elasticity} shows that local participation sensitivity at the observed equilibrium yields producer-side share $\Sigma_D/(\Sigma_S+\Sigma_D)$. The measures often agree qualitatively in linear symmetric environments, but generally differ because the Shapley rule averages marginal contributions over the coalition lattice whereas the local rule takes a derivative at the grand-coalition equilibrium.


\subsection{Remarks}
The Shapley attribution depends on (i) the coalition trading rule and (ii) the harm function. It allocates realized harm without invoking a policy baseline. As emphasized in Section~\ref{sec:contribution_vs_blame}, it is a cooperative-game attribution, not by itself a judgment about causation, fault, or culpability. Section~\ref{sec:elasticity} develops a complementary local sensitivity rule with closed-form expressions in terms of slopes and elasticities at the observed equilibrium.

\section{Local Participation-Sensitivity Allocation via Elasticities}
\label{sec:elasticity}
Shapley values allocate realized harm by averaging marginal contributions across all coalition orderings. A complementary approach keeps the full market fixed and measures local sensitivity: how equilibrium quantity responds to a small proportional scaling of a given agent's supply or demand schedule. This is ordinary comparative statics, not a new equilibrium theorem. Similar gradient-based and path-based allocation ideas appear in the Aumann--Shapley and sensitivity-analysis literatures~\cite{AumannShapley1974,Owen2014,SongNelsonStaum2016}. What is proposed here is the particular perturbation and normalization used for this producer--consumer problem.

\subsection{Proportional participation perturbations}
Introduce participation (scaling) parameters $\lambda_i>0$ for producers and $\mu_j>0$ for consumers:
\begin{equation}
S(p;\lambda)=\sum_{i\in\mathcal{P}} \lambda_i s_i(p),\qquad
D(p;\mu)=\sum_{j\in\mathcal{C}} \mu_j d_j(p),
\end{equation}
with baseline $\lambda_i=\mu_j=1$ for all $i,j$. The (perturbed) market-clearing condition is
\begin{equation}
F(p;\lambda,\mu)\triangleq S(p;\lambda)-D(p;\mu)=0.
\label{eq:F}
\end{equation}
Define log-parameters $\theta_i\triangleq \ln \lambda_i$ and $\psi_j\triangleq \ln \mu_j$. A differential $d\theta_i$ is an equal proportional (percent) perturbation across agents.

Let $p(\lambda,\mu)$ be the clearing price solving \eqref{eq:F}, and define the corresponding cleared quantity
\begin{equation}
Q(\lambda,\mu)\triangleq S(p(\lambda,\mu);\lambda)=D(p(\lambda,\mu);\mu).
\end{equation}
We write $p^\star=p(\mathbf{1},\mathbf{1})$ and $Q^\star=Q(\mathbf{1},\mathbf{1})$ for the baseline equilibrium.

\subsection{Comparative statics at equilibrium}
Assume differentiability and that the implicit-function regularity condition holds at $(p^\star,\mathbf{1},\mathbf{1})$:
\begin{equation}
F_p(p^\star;\mathbf{1},\mathbf{1})=S'(p^\star)-D'(p^\star)>0,
\end{equation}
where $S'(p)=\sum_{i} s_i'(p)$ and $D'(p)=\sum_{j} d_j'(p)\le 0$. Define the positive slope aggregates
\begin{equation}
A \triangleq S'(p^\star)=\sum_{i\in\mathcal{P}} s_i'(p^\star),\qquad
B \triangleq -D'(p^\star)= -\sum_{j\in\mathcal{C}} d_j'(p^\star),
\label{eq:AB}
\end{equation}
so that $F_p=A+B$.

Let each producer's equilibrium quantity be $q_i^\star \triangleq s_i(p^\star)$ and each consumer's equilibrium quantity be $x_j^\star \triangleq d_j(p^\star)$, so $\sum_i q_i^\star=\sum_j x_j^\star=Q^\star$.

\paragraph{Producer sensitivity.}
Differentiate \eqref{eq:F} with respect to $\theta_i$ at baseline. Since $\partial S/\partial \theta_i=\lambda_i s_i(p)$ and $\partial D/\partial \theta_i=0$,
\begin{equation}
\frac{\partial p}{\partial \theta_i}\Big|_\star = -\frac{q_i^\star}{A+B}.
\label{eq:dp_dtheta}
\end{equation}
The derivative of cleared quantity with respect to $\theta_i$ is
\begin{equation}
\frac{\partial Q}{\partial \theta_i}\Big|_\star
=
\underbrace{q_i^\star}_{\text{direct scale}}
+
\underbrace{A\,\frac{\partial p}{\partial \theta_i}\Big|_\star}_{\text{price feedback}}
=
q_i^\star\left(1-\frac{A}{A+B}\right)
=
q_i^\star\,\frac{B}{A+B}.
\label{eq:dQ_dtheta}
\end{equation}

\paragraph{Consumer sensitivity.}
Similarly, for $\psi_j$ we have $\partial D/\partial \psi_j=\mu_j d_j(p)$ and $\partial S/\partial \psi_j=0$, so
\begin{equation}
\frac{\partial p}{\partial \psi_j}\Big|_\star = +\frac{x_j^\star}{A+B},
\end{equation}
and the cleared quantity response is
\begin{equation}
\frac{\partial Q}{\partial \psi_j}\Big|_\star
=
A\,\frac{\partial p}{\partial \psi_j}\Big|_\star
=
x_j^\star\,\frac{A}{A+B}.
\label{eq:dQ_dpsi}
\end{equation}

\paragraph{Normalization.}
Summing \eqref{eq:dQ_dtheta} and \eqref{eq:dQ_dpsi} over all agents yields
\begin{equation}
\sum_{i\in\mathcal{P}}\frac{\partial Q}{\partial \theta_i}\Big|_\star
+
\sum_{j\in\mathcal{C}}\frac{\partial Q}{\partial \psi_j}\Big|_\star
=
\frac{B}{A+B}\sum_i q_i^\star+\frac{A}{A+B}\sum_j x_j^\star
=
Q^\star,
\label{eq:pivotal_sum}
\end{equation}
so the local sensitivity weights aggregate to total quantity. The identity also follows from homogeneity: multiplying every participation parameter $(\lambda,\mu)$ by the same positive scalar leaves the clearing price unchanged and scales $Q$ by that scalar, so $Q$ is homogeneous of degree one in the participation parameters. Euler's theorem then gives \eqref{eq:pivotal_sum}.

\subsection{Elasticity representation}
Define individual point elasticities at equilibrium:
\begin{equation}
\eta_i \triangleq \frac{p^\star}{q_i^\star}\,s_i'(p^\star)\quad(>0),\qquad
\varepsilon_j \triangleq -\frac{p^\star}{x_j^\star}\,d_j'(p^\star)\quad(>0).
\end{equation}
These point elasticities are defined for active agents with positive equilibrium quantities. An inactive agent has zero local scaling weight under the maintained differentiability assumptions and should be handled directly through \eqref{eq:dQ_dtheta} or \eqref{eq:dQ_dpsi}, rather than through a ratio with zero quantity.
Then
\begin{equation}
A=\frac{1}{p^\star}\sum_{i\in\mathcal{P}}\eta_i q_i^\star,\qquad
B=\frac{1}{p^\star}\sum_{j\in\mathcal{C}}\varepsilon_j x_j^\star.
\label{eq:AB_elast}
\end{equation}
Define
\begin{equation}
\Sigma_S \triangleq \sum_{i\in\mathcal{P}}\eta_i q_i^\star,\qquad
\Sigma_D \triangleq \sum_{j\in\mathcal{C}}\varepsilon_j x_j^\star.
\label{eq:SigmaSD}
\end{equation}
Then $\frac{B}{A+B}=\frac{\Sigma_D}{\Sigma_S+\Sigma_D}$ and $\frac{A}{A+B}=\frac{\Sigma_S}{\Sigma_S+\Sigma_D}$.

\subsection{Budget-balanced allocation based on local sensitivities}
Let the local quantity-sensitivity weights be
\begin{equation}
w_{P_i}\triangleq \frac{\partial Q}{\partial \theta_i}\Big|_\star,\qquad
w_{C_j}\triangleq \frac{\partial Q}{\partial \psi_j}\Big|_\star.
\end{equation}
The corresponding derivatives of harm are $H'(Q^\star)w_a$. Their sum is $H'(Q^\star)Q^\star$, which equals $H(Q^\star)$ only for linear harm through the origin. Therefore, for a general nonlinear $H$, the following budget-balanced rule is an \emph{average-damage normalization} of local quantity sensitivities:
\begin{equation}
\beta_a \triangleq H^\star\cdot \frac{w_a}{Q^\star},\qquad a\in\mathcal{N}.
\label{eq:beta_def}
\end{equation}
By \eqref{eq:pivotal_sum}, $\sum_a\beta_a=H^\star$. When $H(Q)=\kappa Q$, this rule also equals the exact local harm derivative $\partial H/\partial\ln\lambda_a$ (or $\partial H/\partial\ln\mu_a$). For nonlinear harm it should not be described as a differential decomposition of the harm level.

Combining \eqref{eq:dQ_dtheta}--\eqref{eq:dQ_dpsi} with \eqref{eq:beta_def}, we obtain explicit formulas:
\begin{align}
\beta_{P_i}
&=
H^\star\cdot \frac{q_i^\star}{Q^\star}\cdot \frac{B}{A+B}
=
H^\star\cdot \frac{q_i^\star}{Q^\star}\cdot \frac{\Sigma_D}{\Sigma_S+\Sigma_D},
\qquad i\in\mathcal{P},
\label{eq:beta_prod}\\
\beta_{C_j}
&=
H^\star\cdot \frac{x_j^\star}{Q^\star}\cdot \frac{A}{A+B}
=
H^\star\cdot \frac{x_j^\star}{Q^\star}\cdot \frac{\Sigma_S}{\Sigma_S+\Sigma_D},
\qquad j\in\mathcal{C}.
\label{eq:beta_cons}
\end{align}
Thus, the producer-side share of this realized-harm allocation is $\Sigma_D/(\Sigma_S+\Sigma_D)$ and the consumer-side share is $\Sigma_S/(\Sigma_S+\Sigma_D)$, with within-side allocations proportional to equilibrium quantities. The cross-side form is the same elasticity logic that appears in standard incidence calculations.

\section{Pigouvian Benchmark and Beneficiary-Pays Allocation of Excess Harm}
The preceding rules allocate the modeled harm at the actual outcome. A different question is how that harm compares with a specified welfare benchmark. Pigouvian analysis supplies a standard benchmark under strong assumptions about welfare, information, and implementability~\cite{Pigou1932,Baumol1972}. It does not, by itself, determine how the benchmark gap should be distributed among agents.

\subsection{Welfare and the efficient quantity}
Assume aggregate demand and supply are invertible (at least locally) so we can define inverse demand and supply:
\begin{equation}
p_D(Q) = D^{-1}(Q),\qquad p_S(Q)=S^{-1}(Q),
\end{equation}
with $p_D'(Q)\le 0$, $p_S'(Q)\ge 0$.

Assume additionally that inverse demand measures marginal willingness to pay, inverse supply measures marginal social opportunity cost exclusive of the external harm, and there are no other distortions. Define utilitarian social welfare (quasi-linear, no income effects) as
\begin{equation}
W(Q)=\int_0^Q p_D(q)\,dq - \int_0^Q p_S(q)\,dq - H(Q).
\label{eq:welfare}
\end{equation}
An interior welfare optimum $Q^{\text{soc}}$ satisfies the first-order condition
\begin{equation}
p_D(Q^{\text{soc}})=p_S(Q^{\text{soc}})+MH(Q^{\text{soc}}).
\label{eq:foceff}
\end{equation}
This is the standard marginal-benefit-equals-marginal-social-cost condition associated with Pigouvian analysis~\cite{Pigou1932,Baumol1972}. Corner solutions require the corresponding inequality conditions.

\subsection{Pigouvian tax and implementation}
Consider a per-unit tax (or fee) $t$ that creates a wedge between consumer and producer prices:
\begin{equation}
p_c - p_p = t.
\end{equation}
With a constant tax $t$, equilibrium must satisfy
\begin{equation}
S(p_p)=D(p_c),\qquad p_c=p_p+t.
\end{equation}
Equivalently, in $p_c$ alone:
\begin{equation}
S(p_c-t)=D(p_c).
\label{eq:wedgeeq}
\end{equation}

A constant per-unit tax
\begin{equation}
t^\star=MH(Q^{\text{soc}})
\label{eq:pigtax}
\end{equation}
is sufficient to implement the interior optimum in this one-good competitive model: at $Q^{\text{soc}}$, market equilibrium satisfies $p_D(Q)-p_S(Q)=t^\star$, which is exactly \eqref{eq:foceff}. A nonlinear schedule $t(Q)=MH(Q)$ can also be written, but it is not required for implementation of the single target quantity. The standard prescription is the marginal external damage evaluated at the efficient output, subject to the usual informational and second-best qualifications~\cite{Baumol1972,Coase1960}.

\subsection{Avoidable harm and the counterfactual gap}
Define the Pigouvian benchmark harm
\begin{equation}
H^{\text{soc}} \triangleq H(Q^{\text{soc}}),
\end{equation}
and the \emph{excess harm relative to the Pigouvian benchmark} as
\begin{equation}
\Delta H \triangleq H(Q^\star)-H(Q^{\text{soc}})=H^\star-H^{\text{soc}} \ge 0.
\label{eq:deltaH}
\end{equation}
The benchmark does not show that $H^{\text{soc}}$ is technologically or morally unavoidable. It is merely the residual harm selected by the stated utilitarian objective and constraints. Under the maintained assumptions, a corrective tax implementing $Q^{\text{soc}}$ removes the gap $\Delta H$.

\subsection{A proposed beneficiary-pays allocation via private gains from underpricing}
Pigouvian theory does not imply an agent-level allocation of $\Delta H$. The following is an additional normative proposal: allocate the benchmark gap in proportion to each agent's private surplus gain from moving from the taxed outcome to the unregulated outcome. This resembles a beneficiary-pays or unjust-enrichment principle rather than a causal theorem.

Let each consumer $j$ have (quasi-linear) indirect surplus at price $p$:
\begin{equation}
CS_j(p)=\int_{p}^{\infty} d_j(x)\,dx,
\label{eq:cs}
\end{equation}
assuming the integral exists (e.g., compact support or sufficiently fast decay). When producer $i$ has a convex variable-cost function $C_i(q)$ with competitive supply $s_i(p)=(C_i')^{-1}(p)$ on the relevant range, producer surplus at producer price $p$ is
\begin{equation}
PS_i(p)=p\,s_i(p)-C_i(s_i(p)).
\label{eq:pscost}
\end{equation}
With the normalization $C_i(0)=0$ and no fixed cost, this is equivalently
\begin{equation}
PS_i(p)=\int_0^{s_i(p)}[p-MC_i(q)]\,dq
=\int_0^p s_i(x)\,dx.
\end{equation}
The first equality, not the sum of these two expressions, is the producer-surplus formula.

Let $(p_c^{\text{soc}},p_p^{\text{soc}})$ denote the consumer and producer prices at the Pigouvian-implemented allocation (with wedge $t(Q^{\text{soc}})$), and let $p^\star$ denote the unregulated market price (where $p_c=p_p=p^\star$).

Define each agent's private gain from underpricing:
\begin{align}
\Delta CS_j &\triangleq CS_j(p^\star)-CS_j(p_c^{\text{soc}}),\\
\Delta PS_i &\triangleq PS_i(p^\star)-PS_i(p_p^{\text{soc}}).
\end{align}
Let total private gain be
\begin{equation}
\Delta \Pi \triangleq \sum_{j\in \mathcal{C}}\Delta CS_j + \sum_{i\in \mathcal{P}}\Delta PS_i.
\end{equation}

\paragraph{Beneficiary-pays allocation.}
Provided $\Delta\Pi>0$, define
\begin{equation}
B_a \triangleq
\Delta H\cdot \frac{\Delta \Pi_a}{\Delta \Pi},
\qquad
\Delta \Pi_a=
\begin{cases}
\Delta PS_a, & a\in \mathcal{P},\\
\Delta CS_a, & a\in \mathcal{C}.
\end{cases}
\label{eq:pigblame}
\end{equation}
Then $\sum_a B_a=\Delta H$. If $\Delta H=0$ or $\Delta\Pi=0$, set $B_a=0$. Under the regular one-good case, $p_c^{\text{soc}}\ge p^\star\ge p_p^{\text{soc}}$, so these gross private-surplus gains are nonnegative. The rule omits tax revenue and its redistribution; it therefore allocates by gross benefit from the missing tax, not by the full welfare effect.

\subsection{Interpretation and variants}
Equation \eqref{eq:pigblame} is a beneficiary-pays allocation: it charges those who receive larger private surplus gains from the uncorrected equilibrium. Benefit is not the same as causation, and the rule says nothing about knowledge, fault, or ability to avoid the harm. Alternative normative rules include polluter-pays, tax-incidence, ability-to-pay, or ability-to-abate principles. The accounting identity is
\begin{equation}
H^\star = H^{\text{soc}} + \Delta H,
\end{equation}
but neither the identity nor Pigouvian efficiency selects a unique agent-level allocation.

\section{Imperfect Information: Ex Ante Expected Attribution}
\label{sec:exante}
So far, curves are treated as known to the analyst. If types are uncertain, any ex post allocation rule can be averaged over a prior. This is a standard expectation operation; it is not by itself a Bayesian equilibrium or a theory of blame. Conditioning on an agent's information describes foreseeable expected attribution, but a culpability analysis would additionally require a causal model, feasible actions, duties, and standards of care~\cite{ChocklerHalpern2004}.

Let each producer have a private type $\theta_i$ determining $s_i(p;\theta_i)$ and each consumer a private type $\eta_j$ determining $d_j(p;\eta_j)$. Let $\omega=(\theta_1,\dots,\theta_N,\eta_1,\dots,\eta_M)$ be the type profile drawn from a common prior $\mathbb{P}$, and let $v_\omega(\cdot)$ be the realized characteristic function induced by $\omega$.

\paragraph{Ex ante expected attribution.}
For any ex post allocation rule $R_a(\omega)$ (e.g., Shapley $\phi_a(v_\omega)$ or the local sensitivity allocation $\beta_a(\omega)$), define ex ante expected attribution as
\begin{equation}
\Phi_a \triangleq \mathbb{E}_\omega\big[R_a(\omega)\big].
\end{equation}

\paragraph{Information-conditioned expected attribution.}
If agent $a$ has information $\mathcal{I}_a$ (e.g., own type and observed price), an information-conditioned expectation can be defined by
\begin{equation}
\Phi_a^{\text{info}} \triangleq \mathbb{E}\big[R_a(\omega)\mid \mathcal{I}_a\big],
\end{equation}
This quantity may be relevant to foreseeability, but it is not itself blame or culpability in the sense of Section~\ref{sec:contribution_vs_blame}.

\section{Symbolic Worked Examples (Two Producers, Three Consumers)}
This section illustrates the mechanics with symbols only, for the case of two producers and three consumers:
\begin{align}
\mathcal{P}&=\{P_1,P_2\},\\
\mathcal{C}&=\{C_1,C_2,C_3\},\\
\mathcal{N}&=\{P_1,P_2,C_1,C_2,C_3\}.
\end{align}

\subsection{Generic coalition quantities}
For any coalition $T\subseteq \mathcal{N}$, define
\begin{equation}
S_T(p)=\sum_{P_i\in T} s_i(p),\qquad
D_T(p)=\sum_{C_j\in T} d_j(p).
\end{equation}
Coalition market clearing yields $p(T)$ such that $S_T(p(T))=D_T(p(T))$ and
\begin{equation}
Q(T)=S_T(p(T))=D_T(p(T)).
\end{equation}
The Shapley game is $v(T)=H(Q(T))$.

\subsection{Shapley contribution: explicit expansion for $P_1$}
With $|\mathcal{N}|=5$, the Shapley value for $P_1$ is
\begin{equation}
\phi_{P_1}
=\mspace{-20mu}
\sum_{T\subseteq \mathcal{N}\setminus\{P_1\}}\mspace{-15mu}
\frac{|T|!\,(4-|T|)!}{5!}\big(H(Q(T\cup\{P_1\}))-H(Q(T))\big).
\end{equation}
Writing $\mathcal{N}\setminus\{P_1\}=\{P_2,C_1,C_2,C_3\}$, we can group terms by $k=|T|$:
\begin{align}
\phi_{P_1}&=\sum_{k=0}^{4} w_k \mspace{-33mu} \sum_{\substack{|T|=k\\T\subseteq \{P_2,C_1,C_2,C_3\} }}\mspace{-33mu}
\big(H(Q(T\cup\{P_1\}))-H(Q(T))\big),\\[1ex]
w_k&\triangleq \frac{k!\,(4-k)!}{5!}.
\end{align}
This expression is fully symbolic; evaluation reduces to computing the relevant coalition equilibrium quantities $Q(T)$.

\subsection{Local elasticity allocation: explicit realized-harm shares}
Let $(p^\star,Q^\star)$ be the unregulated equilibrium. Define equilibrium individual quantities
\begin{align}
q_1^\star&=s_1(p^\star),\; q_2^\star=s_2(p^\star),\\
x_1^\star&=d_1(p^\star),\; x_2^\star=d_2(p^\star),\; x_3^\star=d_3(p^\star),
\end{align}
with $Q^\star=q_1^\star+q_2^\star=x_1^\star+x_2^\star+x_3^\star$.
Define point elasticities
\begin{align}
\eta_1&=\frac{p^\star}{q_1^\star}s_1'(p^\star),\quad
\eta_2=\frac{p^\star}{q_2^\star}s_2'(p^\star),\\
\varepsilon_j&=-\frac{p^\star}{x_j^\star}d_j'(p^\star),\; j=1,2,3.
\end{align}
Let $\Sigma_S=\eta_1 q_1^\star+\eta_2 q_2^\star$ and $\Sigma_D=\varepsilon_1 x_1^\star+\varepsilon_2 x_2^\star+\varepsilon_3 x_3^\star$.
Then the local-sensitivity realized-harm allocation \eqref{eq:beta_prod}--\eqref{eq:beta_cons} gives
\begin{align}
\beta_{P_1} &= H^\star\cdot \frac{q_1^\star}{Q^\star}\cdot \frac{\Sigma_D}{\Sigma_S+\Sigma_D},\\
\beta_{P_2} &= H^\star\cdot \frac{q_2^\star}{Q^\star}\cdot \frac{\Sigma_D}{\Sigma_S+\Sigma_D},\\[2ex]
\beta_{C_1} &= H^\star\cdot \frac{x_1^\star}{Q^\star}\cdot \frac{\Sigma_S}{\Sigma_S+\Sigma_D},\\
\beta_{C_2} &= H^\star\cdot \frac{x_2^\star}{Q^\star}\cdot \frac{\Sigma_S}{\Sigma_S+\Sigma_D},\\
\beta_{C_3} &= H^\star\cdot \frac{x_3^\star}{Q^\star}\cdot \frac{\Sigma_S}{\Sigma_S+\Sigma_D}.
\end{align}
These five shares sum to $H^\star$ by construction.

\subsection{Pigouvian benchmark and contribution in symbols}
Let the unregulated equilibrium satisfy
\begin{equation}
S(p^\star)=D(p^\star),\qquad Q^\star=S(p^\star).
\end{equation}
Let the welfare optimum solve
\begin{equation}
p_D(Q^{\text{soc}})=p_S(Q^{\text{soc}})+MH(Q^{\text{soc}}),
\end{equation}
with corrective tax $t^\star=MH(Q^{\text{soc}})$ and corresponding prices
\begin{align}
p_c^{\text{soc}}&=p_D(Q^{\text{soc}}),\\ p_p^{\text{soc}}&=p_S(Q^{\text{soc}}),\\
p_c^{\text{soc}}-p_p^{\text{soc}}&=MH(Q^{\text{soc}}).
\end{align}
Excess harm relative to the benchmark is
\begin{equation}
\Delta H = H(Q^\star)-H(Q^{\text{soc}}).
\end{equation}
The proposed beneficiary-pays allocation for each actor $a$ is then
\begin{equation}
B_a=\Delta H\cdot \frac{\Delta \Pi_a}{\sum_{b\in \mathcal{N}} \Delta \Pi_b},
\end{equation}
where (using \eqref{eq:cs} and \eqref{eq:pscost} for surplus)
\begin{align}
\Delta CS_j &= CS_j(p^\star)-CS_j(p_c^{\text{soc}}),\qquad j\in\{1,2,3\},\\
\Delta PS_i &= PS_i(p^\star)-PS_i(p_p^{\text{soc}}),\qquad i\in\{1,2\}.
\end{align}

\subsection{Discussion: realized-harm versus benchmark-gap objects}
In the two-producer/three-consumer case, the Shapley and local-sensitivity rules both allocate $H(Q^\star)$ across five players, but they use different operations: Shapley averages marginal contributions across coalition orderings, while the local rule uses marginal influence on $Q^\star$ under proportional curve perturbations. The Pigouvian section first defines $\Delta H$, the difference between the unregulated and internalized outcomes, and then applies a separate beneficiary-pays rule. The three constructions answer different attribution questions (cf.\ Section~\ref{sec:contribution_vs_blame}):
\begin{itemize}
\item \textbf{Shapley:} ``How does the specified coalition cost game allocate realized harm under random-order averaging?''
\item \textbf{Local sensitivity:} ``Who is locally pivotal for the realized quantity, as measured by marginal comparative statics at the observed equilibrium?''
\item \textbf{Pigouvian beneficiary-pays:} ``Who receives the private-surplus gains associated with the benchmark gap?''
\end{itemize}
A complete normative account would require additional principles governing the benchmark residual $H(Q^{\text{soc}})$, the excess $\Delta H$, fault, knowledge, and feasible alternatives. The present paper does not supply that theory.

\section{Conclusion}
This paper proposed three curve-based allocation frameworks for harmful market externalities. The first applies the established Shapley value to a newly specified coalition market-clearing cost game. The second derives a local proportional-participation sensitivity allocation and expresses its producer--consumer split in terms of slopes and elasticities. For nonlinear harm, its budget balance comes from average-damage normalization rather than an exact derivative identity. The third uses the standard Pigouvian welfare benchmark to define an excess-harm gap and then adds a beneficiary-pays allocation based on gross private-surplus gains. The Shapley and Pigouvian foundations are known; the potentially original content lies in the particular game, the symmetric linear finite-sum and limit results, the local elasticity formula in this setting, and the comparative synthesis. None of the three rules is, without additional causal and normative assumptions, a complete measure of actual causation, legal liability, or moral blame.

Several extensions are immediate: incorporating heterogeneous emissions intensity and avoidable-technology choice, adding abatement choices, treating uncertainty formally via Bayesian equilibria, evaluating alternative welfare criteria (rights-based constraints, caps, distributional objectives) that change the definition of the ``acceptable'' benchmark harm, and developing a fault-sensitive layer on top of the model-based allocations presented here.

\section*{Acknowledgment}
The author wishes to acknowledge the use of large language models (Sonnet 4.7 and GPT 5.5) for assistance in improving the language and clarity of the manuscript, and for assisting with code and algebraic checking. The author remains responsible for the claims, derivations, citations, and any remaining errors.

\begin{thebibliography}{99}

\bibitem{Shapley1953}
L.~S. Shapley, ``A value for $n$-person games,'' in \emph{Contributions to the Theory of Games II}, H.~W. Kuhn and A.~W. Tucker, Eds. Princeton, NJ, USA: Princeton Univ. Press, 1953, pp. 307--317.

\bibitem{AumannShapley1974}
R.~J. Aumann and L.~S. Shapley, \emph{Values of Non-Atomic Games}. Princeton, NJ, USA: Princeton Univ. Press, 1974.

\bibitem{ShapleyShubik1969}
L.~S. Shapley and M.~Shubik, ``On market games,'' \emph{Journal of Economic Theory}, vol.~1, no.~1, pp. 9--25, 1969.

\bibitem{ThrallLucas1963}
R.~M. Thrall and W.~F. Lucas, ``$n$-person games in partition function form,'' \emph{Naval Research Logistics Quarterly}, vol.~10, no.~4, pp. 281--298, 1963.

\bibitem{Pigou1932}
A.~C. Pigou, \emph{The Economics of Welfare}, 4th ed. London, U.K.: Macmillan, 1932, Part II, ch.~IX, ``Divergences between marginal social net product and marginal private net product.''

\bibitem{Coase1960}
R.~H. Coase, ``The problem of social cost,'' \emph{Journal of Law and Economics}, vol.~3, pp. 1--44, 1960.

\bibitem{Baumol1972}
W.~J. Baumol, ``On taxation and the control of externalities,'' \emph{American Economic Review}, vol.~62, no.~3, pp. 307--322, 1972.

\bibitem{BilleraHeath1982}
L.~J. Billera and D.~C. Heath, ``Allocation of shared costs: A set of axioms yielding a unique procedure,'' \emph{Mathematics of Operations Research}, vol.~7, no.~1, pp. 32--39, 1982.

\bibitem{Owen2014}
A.~B. Owen, ``Sobol' indices and Shapley value,'' \emph{SIAM/ASA Journal on Uncertainty Quantification}, vol.~2, no.~1, pp. 245--251, 2014.

\bibitem{SongNelsonStaum2016}
E.~Song, B.~L. Nelson, and J.~Staum, ``Shapley effects for global sensitivity analysis: Theory and computation,'' \emph{SIAM/ASA Journal on Uncertainty Quantification}, vol.~4, no.~1, pp. 1060--1083, 2016.

\bibitem{HalpernPearl2005}
J.~Y. Halpern and J.~Pearl, ``Causes and explanations: A structural-model approach. Part I: Causes,'' \emph{British Journal for the Philosophy of Science}, vol.~56, no.~4, pp. 843--887, 2005.

\bibitem{ChocklerHalpern2004}
H.~Chockler and J.~Y. Halpern, ``Responsibility and blame: A structural-model approach,'' \emph{Journal of Artificial Intelligence Research}, vol.~22, pp. 93--115, 2004.

\bibitem{HeskesEtAl2020}
T.~Heskes, E.~Sijben, I.~G. Bucur, and T.~Claassen, ``Causal Shapley values: Exploiting causal knowledge to explain individual predictions of complex models,'' in \emph{Advances in Neural Information Processing Systems 33}, 2020, pp. 4778--4789.

\bibitem{XieChen2023}
R.~Xie and Y.~Chen, ``Real-time bidding strategy of energy storage in an energy market with carbon emission allocation based on Aumann--Shapley prices,'' arXiv:2307.00207, 2023.

\bibitem{FengEtAl2025}
Y.~Feng, T.~Sun, Y.~Meng, X.~Yang, and D.~Feng, ``Deep learning-accelerated Shapley value for fair allocation in power systems: The case of carbon emission responsibility,'' arXiv:2511.01229, 2025.

\end{thebibliography}

\end{document}
